%!TEX root = main.tex

\section{Braided Operads}
    Braided operads were first mentioned by Fiedorowicz in his paper \cite{EO} on construction of $\E_n$ operads. Let $A$ be an operad in the  category of compactly generated Hausdorff spaces. The first theorem states:

    \begin{thm}[Recognition principles for $E_n$ operads]
        $A$ is an $\E_2$ operad iff $A(k)$ is connected and the collection of universal covering spaces $\{\Tilde{A}(k)\}_{k \in \N}$ forms a $B_{\infty}$ operad, i.e.:
        \begin{enumerate}[(1)]
            \item For all $k \in \N$ $\Tilde{A}(k)$ is contractible.
            \item The braid group  $B_k$ acts freely on $\Tilde{A}(k)$ for each $k \in \N$.
        \end{enumerate}
    \end{thm}

    In the proof, it is stated that "the definition of operads can be reformulated using actions by braid groups in place of actions by symmetric groups". As braided operads are a useful tool, e.g. recognizing $\E_n$ operads, we want to gain a deeper understanding of braided operads. Further, we have already seen that colored operads correspond to fibrations over finite pointed sets $\Fin$ as well as colored planar operads correspond to fibrations over $\deltast$. We are hoping to find a suitable category such that colored braided operads correspond to fibrations over this category. 
    
    % To do so, we first need to develop the language for it. Meaning we will need to develop a concept of braided colored operads which we will work towards in this section. Doing so we define action operads which lead us to group operads and group colored operads. Braided colored operads are then a special case of these where the group are in fact the braid groups.

    It might be interesting to try to define braided operads in braided monoidal categories rather than symmetric monoidal categories potentially yielding interesting examples. However, it is not straight forward how to do this as shuffling objects would not be a canonical isomorphism. So for now, we will stay in a symmetric monoidal category $\C$.

\begin{defin}
    A {\em braided operad} $\O$ in $\C$ is a planar operad as defined in \ref{planaroperad} with an action of $B_n$ on $\O(n)$ for all $n \in \N$ together with the following equivariance condition: 
           For $\sigma \in B_k$, let $\sigma(j_1, \cdots, j_k) \in B_j$ denote the element that $\sigma$-braids $k$ blocks with each block having $j_l$ strands for $1 \leq l \leq k$. Then the following diagram commutes:
                \begin{center}
                \begin{tikzcd}
                    \O(k) \otimes \O(j_1) \otimes \cdots \otimes \O(j_k) \arrow[d, "\gamma"'] \arrow[rr, "\sigma \otimes \bar{\sigma}^{-1}"] &  & \O(k) \otimes \O(j_{\bar{\sigma}(1)}) \otimes \cdots \otimes \O(j_{\bar{\sigma}(k)}) \arrow[d, "\gamma"] \\
                    \O(j) \arrow[rr, "\sigma(j_1 \cdots j_k)"']                                                                                        &  & O(j),                                                                   
                \end{tikzcd}
                \end{center}      
            where $\bar{\sigma}$ denotes the underlying permutation of $\sigma$.
            Furthermore, let $\tau_l \in B_{j_l}$ for $1 \leq l \leq k$ and $\tau_1 \oplus \cdots \oplus \tau_k$ denotes the sum of braidings on $k$ blocks with each block having $j_l$ strands. Then the following diagram commutes
                
                \begin{center}
                \begin{tikzcd}
                    \O(k) \otimes \O(j_1) \otimes \cdots \otimes \O(j_k) \arrow[d, "\gamma"'] \arrow[rr, "\id \otimes \tau_1 \otimes \cdots \otimes \tau_k"] &  & \O(k) \otimes \O(j_{\sigma(1)}) \otimes \cdots \otimes \O(j_{\sigma(k)}) \arrow[d, "\gamma"] \\
                    \O(j) \arrow[rr, "\tau_1 \oplus \cdots \oplus \tau_k"']                                                                                &  & \O(j)                                                                                     
                \end{tikzcd}
                \end{center}
\end{defin}

\begin{rem}
    Similarily to symmetric operads, we can define algebras and modules over braided operads by changing the equivariance condition:

    Let $\O$ be a braided operad and $A$ an algebra over the underlying planar operad of $\O$. For all $j \in \N$, $\sigma \in B_j$ and $\bar{\sigma} \in S_j$ being the underlying permutation of $\sigma$:
    \begin{center}
        \begin{tikzcd}
            O(j) \otimes A^j \arrow[rr, "\sigma \otimes \bar{\sigma}^{-1}"] \arrow[rd, "\theta"'] &   & O(j) \otimes A^j \arrow[ld, "\theta"] \\
            & A &     
        \end{tikzcd}
    \end{center}
    Then, we call $A$ an algebra over the braided operad $\O$.

    Further, let $A$ be an algebra over the braided operad $\O$ and $M$ an object in $\C$ with
    \[
        \omega \colon O(j) \otimes A^{j-1} \otimes M \rightarrow M 
    \]
    satisfying associativity and unitality conditions from \ref{module}.
    Then we call $M$ an $A$-module if the following equvariance condition holds for all $j \in \N$, $\sigma \in B_{j-1} \subset B_j$ and $\bar{\sigma} \in S_{j-1}$ being the underlying permutation of $\sigma$:

    \begin{center}
                \begin{tikzcd}
                    O(j) \otimes A^{j-1} \otimes M \arrow[rd, "\omega"'] \arrow[rr, "\sigma \otimes \bar{\sigma}^{-1} \otimes \id"] &   & O(j) \otimes A^{j-1} \otimes M \arrow[ld, "\omega"] \\
                     & M &                                                    
                \end{tikzcd}
                \end{center}
\end{rem}

\begin{eg}
    \begin{enumerate}
        \item Again, let us consider the operad in $\Set$ defined by the singleton set in each degree $n \in \N$. The braid group acts trivially. When we consider algebras over this operad we will again obtain commutative algebras.
        \item We have already encountered an example of planar operad in \ref{braidop} which can be lifted to a braided operad---the braid operad $\B$. The braid group $B_n$ acts by composition. This group action is equivariant.
    \end{enumerate}
    
\end{eg}

\subsection{Action Operad}
For this section, we remain in the symmetric monoidal category of small sets equipped with the direct product as the monoidal product. We will introduce an action operad that acts on planar operads and we will see how this relates to braided operads following \cite{yau}.

\begin{defin}
    An {\em action operad} is a tuple
    \[
        (\{G(n)\}_{n \in \N}, \gamma^G, 1^G, \omega)
    \]
    satisfying the following conditions:
    \begin{enumerate}[(i)]
        \item For all $n \in \N$, $G(n)$ is a group with unit $\id_n$.
        \item There is a group homomorphism $\omega \colon G(n) \to S_n$ for all $n$ which we call augmentation. For each $g \in G(n)$ we call $\omega(g)$ the underlying permutation of $g$.
        \item $(\{G(n)\}_{n \in \N}, \gamma^G, 1^G)$ is a one-colored planar operad in $\Set$ which we denote by $G$. It fulfills the following conditions:
        \begin{enumerate}[(a)]
            \item The augmentation
            \[
            \omega \colon G \to \S
            \]
            is a morphism of planar operads in $\Set$ where $\S$ denotes the symmetry operad from \ref{egop}.
            \item For $n \geq 1$, let $\sigma, \rho \in G(n)$ and for $k_j \in \N$, $1 \leq j \leq n$ let $\tau_j, \tau_j' \in G(k_j)$ . Then the following equation holds in $G(k)$ for $ \sum_{1 \leq j \leq n} k_j = k$:
            \[
                \gamma^G(\sigma \rho, \tau_1 \tau_1', \dots, \tau_n \tau_n') = \gamma^G(\sigma, \tau_{\bar{\rho}^{-1}(1)}, \dots, \tau_{\bar{\rho}^{-1}(n)}) \cdot \gamma^G(\rho, \tau_1', \dots, \tau_n')
            \]
            where $\bar{\rho} = \omega(\rho) \in S_n$.
        \end{enumerate}
    \end{enumerate}
\end{defin}
We call the condition (iii)(b) the action operad equation.

\begin{egs}
    \begin{enumerate}
        \item The \textit{planar group operad} $P$ is the action operad where $P(n) = \{\ast\}$ is the trivial group for all $n \in \N$. 
        \item The symmetry operad $\S$ from \ref{egop} equipped with the identity augmentation forms an action operad.
        \item Another example which we have already seen is the braid operad which we equip with the projection to the symmetric group as the augmentation.
    \end{enumerate}
\end{egs}

\begin{prop}
    The braid operad $\B$ from \ref{braidop} equipped with the projection 
    \[
        \pi \colon \B \to \S
    \]
    forms an action operad
    \[
        B = \left( \{B_n\}_{n \in \N}, \gamma, \id_1 \in \B_1, \pi \right)
    \]
\end{prop}

\begin{proof}
    We have already seen that $\B$ forms a one-colored planar operad in $\Set$. Hence, it remains to show that $\pi$ is a morphism of planar operads.
    We note that $\pi(\id_1)$ is the trivial permutation and $\pi$ preserves the operadic composition as it perserves the group multiplication. The action operad equation follows from \ref{braidfacts}.
\end{proof}

\subsection{Outlook}
    We want to define braided operads as fibrations over a suitable category. A first intuition led us to the assumption that colored braided operads might fiber over the crossed simplicial braid group as defined in \cite{Dyckerhoff}.
    
\begin{defin}\label{crossed_braid} The {\em crossed simplicial braid group} is a category $\Delta B$ equipped with an embedding
	$\iota:\Delta \to \Delta B$ such that:
	\begin{enumerate}
		\item the functor $\iota$ is bijective on objects,
		\item any morphism $u: \iota[m] \to \iota[n]$ in $\Delta B$ can be uniquely expressed as 
			a composition $\iota(\phi) \circ b$ where $\phi: [m] \to [n]$ is a morphism in
			$\Delta$ and $b \in B_m$ is acting on $\iota[m]$ in $\Delta B$.
	\end{enumerate}
    Note that we sometimes refer to $\iota[m]$ by writing $[m]$ instead.
\end{defin}

\begin{rem}
    To understand composition in the category $\Delta B$ let us consider two morphisms
    \[
    f \colon [m] \to [n] \quad \text{and} \quad g \colon [n] \to [l]
    \]
    with the factorization $f = \phi_f \circ b_f$ and $g = \phi_g \circ b_g$. 
    We define for every $j \in [n]$ an $f$-pencil $p_j$ as a collection $[j, i_1, \dots i_s]$ with $\phi_f^{-1}(j) = \{i_1 < \dots < i_s\}$. We call $s$ the width of the pencil. Note that we say it is an empty pencil $p_j$ if $\phi_f^{-1}(j) = \emptyset$.
    
    Then we obtain a braid on $m$ strands in the following way:
    First, we consider $p_1, \dots, p_n$ as $n$ strand such that we can braid by $b_g$. Then, we forget the empty pencil strands; in other words we restrict the action of $b_g$ to the subset $J \subset [n]$ such that $p_j$ is not empty for all $j \in J$. However, for notation reason and without loss of generality, we assume that there are no empty pencils.    
    We obtain a braiding on $m$ strands if we consider the action of $b_g$ on the pencils as a block braiding on $m$ strand. Each pencil corresponds to a block having the length of the pencils width. The blocks are ordered by $[n]$.

    Geometrically this can be thought of as braiding the pencils. Then, as the order within the pencil remains unchanged, we can cut a pencil $[j, j_1, \dots, j_s]$ into its $s$ underlying strands. The resulting braid corresponds to $b_g(s_1, \dots, s_n) \circ b_f$ where $s_j$ is the width of the pencil $p_j$ for $j \in [n]$.
    Thus, we define
    \[
        g \circ f \coloneqq \phi_g \circ (b_g(s_1, \dots, s_n) \circ b_f).
    \]
\end{rem}

We can define a new colored operad with the crossed simplicial braid group as its indexing sets. However, this causes some issues. A braided monoidal category is supposed to be an example of the new colored operad, meaning we require a morphism in $\Delta B$, inducing the braiding in the braided monoidal category. Thus, the following diagram would have to commute in $\Delta B$:

\begin{center}
    \begin{tikzcd}
        {[1]} \arrow[r, "\tau^1_0"] \arrow[d, "\sigma_0"'] & {[0]} \\
        {[1]} \arrow[ru, "\tau^1_0"']                      &      
    \end{tikzcd}
\end{center}

which is not the case. 

The next step will be to investigate whether we can identify further relations to define a suitable category related to $\Delta B$.

\todo{add potential fix and explain the contradiction of braid action on cartesian product}

